% AY2023 制御工学 第2問〔II〕（転記）
% 出典: 03_制御工学/original/03_制御工学/2023/2023se.pdf
% 要約: 1/M, 1/s, 1/s, α の直列に D（速度帰還）・K（位置帰還）を持つフィードバック系。
%       伝達関数・ゲイン位相・ボード線図・α 変更時の比較考察。構造は AY2014/2018 第2問と同型。

\section*{第2問〔II〕}

下図の制御システムに関する次の問い (1)〜(4) に答えよ. なお, 入力を $U(s)$, 出力を $Y(s)$
とする. $M$, $D$, $K$, $\alpha$ は係数である.

\begin{center}
\begin{tikzpicture}[>=Latex]
  \node (in)   at (0,0)   {$U(s)$};
  \node (sum1) [sum] at (1.3,0) {};
  \node (sum2) [sum] at (2.7,0) {};
  \node (Minv) [block] at (4.2,0) {$\dfrac{1}{M}$};
  \node (int1) [block] at (5.9,0) {$\dfrac{1}{s}$};
  \coordinate (nodeA) at (7.0,0);
  \node (int2) [block] at (8.1,0) {$\dfrac{1}{s}$};
  \coordinate (nodeB) at (9.2,0);
  \node (al)   [block] at (10.2,0) {$\alpha$};
  \node (out)  at (11.6,0) {$Y(s)$};
  \node (D)    [block] at (5.2,-1.5) {$D$};
  \node (K)    [block] at (5.2,-3.0) {$K$};

  \draw[->] (in)   -- (sum1);
  \draw[->] (sum1) -- (sum2);
  \draw[->] (sum2) -- (Minv);
  \draw[->] (Minv) -- (int1);
  \draw[->] (int1) -- (int2);
  \draw[->] (int2) -- (al);
  \draw[->] (al)   -- (out);
  \fill (nodeA) circle (1.4pt);
  \fill (nodeB) circle (1.4pt);

  \draw (nodeA) -- (7.0,-1.5) -- (D.east);
  \draw[->] (D.west) -| (sum2);
  \draw (nodeB) -- (9.2,-3.0) -- (K.east);
  \draw[->] (K.west) -| (sum1);

  \node at ($(sum1.north west)+(0.02,0.10)$) {\scriptsize $+$};
  \node at ($(sum1.south west)+(0.02,-0.10)$) {\scriptsize $-$};
  \node at ($(sum2.north west)+(0.02,0.10)$) {\scriptsize $+$};
  \node at ($(sum2.south west)+(0.02,-0.10)$) {\scriptsize $-$};
\end{tikzpicture}
\end{center}

\begin{enumerate}[label=(\arabic*)]
  \item 図で表されるシステムの伝達関数 $G(s)$ を求めよ.

  \item 各係数を $M=5$, $D=5.5$, $K=0.5$, $\alpha=5$ とするとき,
        $G(s)$ のゲインならびに位相を求めよ.

  \item 問い (2) で求めたシステムのボード線図のうち, ゲイン線図を漸近線によって描け.
        折れ点等の特徴点を明示し, ゲイン線図の縦軸・横軸には適宜数値を記入すること.

  \item 次に, 問い (2) で求めたシステムにおいて, $\alpha=0.5$ となる場合を考える. このときの
        ゲイン線図を問い (3) で描いたゲイン線図上に破線で描け. また, $\alpha=5$ ならびに
        $\alpha=0.5$ となるそれぞれのシステムに対して単位ステップ入力を与えた時の定常値を
        求めたうえで, 得られた解やゲイン線図などから, 2 つのシステムに関する応答特性の差異に
        ついて考察せよ.
\end{enumerate}
